%======================================================================
% CAPÍTULO 3: MARCO LEGAL
%======================================================================
\chapter{Marco Legal}

\section{Ley de Infogobierno}

La Ley de Infogobierno establece normas para el uso de tecnolog\'{\i}as de informaci\'on en el sector p\'ublico, promoviendo la Soberan\'{\i}a Tecnol\'ogica y el uso de Software Libre.

\begin{table}[H]
\centering
\caption{Cumplimiento con la Ley de Infogobierno}
\begin{tabular}{|l|l|}
\hline
\textbf{Requisito} & \textbf{Implementaci\'on} \\ \hline
Software Libre & Todas las tecnolog\'{\i}as son c\'odigo abierto \\ \hline
Est\'andares Abiertos & REST, gRPC, JSON, Protocol Buffers \\ \hline
Interoperabilidad & APIs documentadas, adaptadores \\ \hline
Seguridad & Cifrado AES-256, JWT, MFA, TLS 1.3 \\ \hline
\end{tabular}
\end{table}

\section{Ley Especial contra los Delitos Inform\'aticos}

Esta ley establece conductas penalizadas relacionadas con el uso indebido de sistemas inform\'aticos. Sandra Ecosystem implementa controles que ayudan a prevenir y detectar estas conductas:

\begin{itemize}
\item Acceso no autorizado: Autenticaci\'on multifactor
\item Interceptaci\'on de datos: Cifrado en tránsito y en reposo
\item Falsificaci\'on de datos: Firmas digitales y logs de auditoría
\item Destrucci\'on de datos: Sistemas de backup y redundancia
\end{itemize}

\section{Ley sobre Mensajes de Datos}

Establece el marco legal para las firmas electr\'onicas y los mensajes de datos. El ecosistema implementa:
\begin{itemize}
\item Firmas digitales reconocidas
\item Sellos de tiempo
\item Certificados digitales
\end{itemize}

\section{Est\'andares ISO}

\begin{caja}{Est\'andares ISO Aplicables}
\begin{itemize}
\item \textbf{ISO/IEC 27001}: Sistema de Gesti\'on de Seguridad de la Informaci\'on
\item \textbf{ISO 22301}: Sistema de Gesti\'on de Continuidad del Negocio
\item \textbf{ISO/IEC 27002}: Controles de seguridad
\item \textbf{NIST SP 800-53}: Controles de seguridad y privacidad
\end{itemize}
\end{caja}

\section{Estrategia y Defensa: Licenciamiento Soberano}

La elección de licencias en el Proyecto Sandra no es solo legal, es \textbf{estratégica}. Priorizamos licencias que garantizan la autonomía tecnológica y la protección de la propiedad intelectual en entornos críticos orientados a la defensa y la soberanía del Estado.

\subsection{Privatizaci\'on de C\'odigo Derivado}
Las licencias permisivas permiten que innovaciones basadas en código abierto puedan ser "privatizadas" y clasificadas, un requerimiento fundamental para el sector defensa y la seguridad nacional.

\subsection{Cumplimiento y Clasificaci\'on}
El ecosistema hace uso estratégico de licencias permisivas como Apache 2.0 y MIT. Esto se hace para permitir que módulos especiales, como algoritmos de guiado o componentes altamente sensibles, sean eventualmente cerrados y protegidos bajo normativas de seguridad nacional, manteniendo abierto solo el ecosistema general.

\subsection{An\'alisis de Impacto Estrat\'egico}
\begin{table}[H]
\centering
\caption{Recomendaciones Estratégicas de Licenciamiento}
\begin{tabular}{|l|p{8cm}|}
\hline
\textbf{Licencia} & \textbf{Recomendación} \\ \hline
\textbf{Apache 2.0 / MIT} & \textbf{Recomendado:} Permite clasificar módulos críticos en un esquema de código propietario cerrado. \\ \hline
\textbf{GPL v3 / AGPL} & \textbf{Restringido:} Su persistencia en el licenciamiento ("copyleft" fuerte) prohíbe el cierre estratégico de los códigos derivados. \\ \hline
\end{tabular}
\end{table}

\notacientifica{<<Si desarrollas un módulo de guiado basado en un algoritmo abierto, una licencia como Apache 2.0 permite que el software final sea cerrado y clasificado sin comprometer el cumplimiento legal.>>}

\newpage
